\documentclass[11pt,a4paper]{article}
\usepackage[margin=25mm]{geometry}
\usepackage{amsmath,amssymb,amsthm,hyperref}
\setlength{\emergencystretch}{2em}
\newtheorem{theorem}{Theorem}
\newtheorem{lemma}[theorem]{Lemma}
\newtheorem{proposition}[theorem]{Proposition}
\newcommand{\gen}[1]{\langle#1\rangle}
\newcommand{\F}{\mathbf F}
\newcommand{\Sp}{\operatorname{Sp}}
\newcommand{\Tr}{\operatorname{Tr}}
\title{Endpoint variation in orthogonal symmetry markings}
\author{Human Compression project: L031 corrected recovery}
\date{October 2026}
\begin{document}\maketitle
We rework the positive characteristic-two construction of
Brooksbank--Ferrara--Leemans [BFL]. Full orders are imposed on the
interior bonds of a path; its two endpoint coefficients may vary.
This supplies both orthogonal types without an unstated choice between
primitive torus parameters. Every group below acts on its actual
ambient space; restriction to a degenerate subspace is not assumed
faithful. The complete source contract is supplied by the companion source-interface
paper and its explicitly acquired low-dimensional and binary proofs.

\section{Classical inputs and support groups}
Let $F=\F_q$, $q=2^k\geq4$. A quadratic space will be called
\emph{nonsingular} here when its polar form $B$ is nondegenerate,
or its radical is one-dimensional and has nonzero quadratic value.
For $Q(v)\ne0$ write
\[
 s_v(x)=x+\frac{B(x,v)}{Q(v)}v,\qquad
 \Sigma_U=\gen{s_v:v\in U,\ Q(v)\ne0}.
\]
An orthogonal subspace has polar radical of dimension at most one.
The following standard classical inputs remain explicit:
(i) in a nondegenerate even-dimensional quadratic space over $F$,
the full orthogonal group is generated by its symmetries;
(ii) the Dickson character of this group is
$g\mapsto\operatorname{rank}(g-1)\pmod2$;
(iii) the nonsplit dihedral subgroup of order $2(q+1)$ is maximal
in $\mathrm{SL}_2(q)$ for $q\geq4$; and
(iv) the Arf trace distinguishes the two nondegenerate even-dimensional
quadratic forms. These are classical group and quadratic-form facts,
not invocations of the construction theorem being recovered.
Input (i) is the symmetry-generation theorem cited by [BFL, \S3]
from Taylor, Theorem 11.39. Input (ii) is Taylor, Theorem 11.43 (printed p.160).
Input (iii) is the rank-one subgroup classification used in
[BFL, Lemma 3.4] and stated in [LMO, Proposition 2.3, p.5]: a subgroup containing an element of order $q+1$
cannot lie in a Borel, split torus normalizer, or proper subfield
group, and the remaining nonsplit torus normalizer is the given one.
Their full proofs remain external acquisition costs.

We will use two elementary consequences. Symplectic groups over $F$
are generated by their transvections
$t_{u,a}(x)=x+aB(x,u)u$, $a\ne0$.
Indeed such maps carry any nonzero vector to any other: use the
transvection in their sum if their polar product is nonzero, and an
intermediate vector having nonzero polar products with both otherwise.
After fixing a vector $e$, two possible partners $f,f'$ with
$B(e,f)=B(e,f')=1$ are joined by transvections with axes in $e^\perp$.
If necessary use $f''=f'+te$ where
$t\notin\{0,B(f,f')\}$; the two polar products are then nonzero.
Having fixed a hyperbolic pair, induct on its perpendicular complement.
Also $\mathrm{SL}_2(q)$ is perfect: its upper and lower elementary
unipotents generate it, and conjugation by
$\operatorname{diag}(a,a^{-1})$, $a^2\ne1$, expresses every such
unipotent as a commutator. Both assertions require no finite census.

\begin{lemma}[Support characterization]\label{support}
Let $V$ be nonsingular as defined above. For an orthogonal subspace
$U\leq V$ of dimension at least two put
\[
 K(U)=\{g\in O(V):\operatorname{im}(g-1)\subseteq U\}.
\]
Then $K(U)=\Sigma_U$. For a nonsingular one-space
$U=\gen v$ the same equality is $K(U)=\gen{s_v}$.
Consequently if $U,W,U\cap W$ are orthogonal subspaces of
dimension at least two, then
$\Sigma_U\cap\Sigma_W=\Sigma_{U\cap W}$.
The conclusion also holds when $U\cap W$ is a nonsingular one-space.
\end{lemma}
\begin{proof}
Every product of symmetries in $U$ has support in $U$.
Conversely, an isometry with support in $U$ fixes $U^\perp$
pointwise. To see this, for $y\in U^\perp$,
$B(x,(g^{-1}-1)y)=B((g-1)x,y)=0$ for every $x$.
Thus $(g^{-1}-1)y$ is in the polar radical of $V$.
It is zero when the radical is zero; otherwise preservation of $Q$
and the nonzero quadratic value of its generator force it to be zero.
In particular $K(U)$ preserves $U$ and fixes its radical.

If $U$ is nondegenerate, $V=U\perp U^\perp$ and
$K(U)$ is its full orthogonal group extended trivially; input (i)
applies. Suppose next that $U=M\perp\gen z$, with $M$
nondegenerate and $Q(z)=0$. This $z$ is not radical in $V$.
Choose $w\in M^\perp$ with $B(z,w)=1$, and replace $w$ by
$w+Q(w)z$ to make it singular. On $M\perp\gen{z,w}$,
the kernel of the induced map $K(U)\longrightarrow O(M)$ consists
exactly of the maps
\[
 T_u(x)=x+B(x,z)u+\bigl(B(x,u)+Q(u)B(x,z)\bigr)z,
 \qquad u\in M.                                      \tag{1}
\]
On its perpendicular complement these maps are the identity.
Indeed the action on $M$ must be $x\mapsto x+B(x,u)z$;
preserving the pairing with $w$ determines its $M$-component $u$,
and preserving $Q(w)$ determines the remaining coefficient $Q(u)$.
Direct substitution gives $T_uT_v=T_{u+v}$.
For $v\in M$ nonsingular and $a\in F$ one has
\[
                  s_v s_{v+az}=T_{av/Q(v)}.            \tag{2}
\]
Thus every scalar multiple of every nonsingular $v$ occurs as a
translation parameter. Such vectors span $M$: in a hyperbolic plane
one may replace its singular basis by $e+f,e+af$ with
$a\notin\{0,1\}$, and an asingular plane already has a nonsingular
basis. Input (i), equations (1)--(2), and the trivial extensions
of $O(M)$ now give $K(U)=\Sigma_U$.

If instead $Q(z)\ne0$, the induced map
$K(U)\longrightarrow\Sp(U/\gen z)$ is onto already on
$\Sigma_U$. In $U=M\perp\gen z$, the transvection with axis
$u\in M$ and coefficient $a\ne0$ is induced by
$s_{u+tz}$, where
$Q(u)+t^2Q(z)=a^{-1}$; the square map on $F$ is bijective.
Use the elementary transvection generation proved above.
An element of the kernel fixes $U$ pointwise, since a nonzero
radical shear inside $U$ would change its quadratic value.
If $z$ is radical in $V$, this kernel is trivial. Otherwise an
element fixing $U$ pointwise and having support in $U$ has the
form $x\mapsto x+cB(x,z)z$. Preservation of $Q$ gives
$c^2Q(z)+c=0$, so the kernel is exactly $\{1,s_z\}$,
also contained in $\Sigma_U$.

For a nonsingular one-space, write a supported map as
$g(x)=x+f(x)v$. Quadratic preservation says
$f(x)(Q(v)f(x)+B(x,v))=0$ for every $x$.
A vector space is not the union of two proper linear subspaces,
so either $f=0$ or $f=B(-,v)/Q(v)$.
Finally intersection of two support conditions is precisely support
in the intersection. This proves every asserted equality.
\end{proof}

\begin{lemma}[An ambient three-space]\label{three}
Let $T\leq V$ be an orthogonal three-space, let $L\leq T$
be asingular of dimension two, and let $x\in T\setminus L$
be nonsingular and nonradical in $T$. Then
\[
                    \gen{\Sigma_L,s_x}=\Sigma_T.
\]
\end{lemma}
\begin{proof}
Write $T=L\perp\gen z$. If $Q(z)=0$, write $x=u+tz$
with $u\ne0$, $t\ne0$. Equation (2) supplies a nonzero
translation in $K(T)=L^+\rtimes O^-(L)$.
The cyclic torus of order $q+1$ acts irreducibly on $L$ even
as an $\F_2$-space: an element of that order in $\F_{q^2}$
has degree $2k$ over $\F_2$, since a proper divisor of $2k$
is at most $k$ and $q+1>2^k-1$. Its orbit therefore spans
every translation parameter over $\F_2$. This gives all of $K(T)$.

Suppose $Q(z)\ne0$. The projection of $\Sigma_L$ to
$\Sp(T/\gen z)=\mathrm{SL}_2(q)$ is its nonsplit dihedral
group. The projection of $s_x$, for $x=u+tz$, has axis $u$
and transvection coefficient $(Q(u)+t^2Q(z))^{-1}$.
This differs from the coefficient $Q(u)^{-1}$ of the unique
dihedral involution on that axis. Hence it is outside that dihedral
group, and input (iii) gives the full $\mathrm{SL}_2(q)$ image.

If $z$ is radical in $V$, projection is injective and we are done.
Otherwise the kernel is $\gen{s_z}$ by Lemma~\ref{support}.
The Dickson character on the nondegenerate even space obtained by
adjoining a partner to $z$ restricts to $K(T)$ and takes value one
on $s_z$. Thus $K(T)$ is the direct product of its character kernel,
isomorphic to $\mathrm{SL}_2(q)$, and $\gen{s_z}$.
A subgroup projecting bijectively to this perfect group would have
trivial Dickson character. Our subgroup contains symmetries, of
character one, so its projection kernel is nontrivial and it equals
$K(T)=\Sigma_T$.
\end{proof}

\begin{lemma}[Hyperplane saturation]\label{saturation}
Let $W$ be an orthogonal hyperplane of a nonsingular quadratic space
$V$, with $\dim V\geq4$. Let $L$ be an asingular two-space not
contained in $W$. Then $\gen{\Sigma_W,\Sigma_L}=\Sigma_V$.
\end{lemma}
\begin{proof}
For a nonsingular $y$ outside $W\cup L$, put $T=\gen{L,y}$.
This is an orthogonal three-space. Its intersection $M=T\cap W$
is a two-space meeting $L$ in a nonsingular point. Among the
$q+1\geq5$ projective points of $M$, at most two are singular,
at most one lies in $L$, and at most one is radical in $T$.
A remaining point supplies a nonsingular nonradical $x\in M\setminus L$.
Lemma~\ref{three} now gives $\Sigma_T$, including $s_y$.
All nonsingular roots in $W\cup L$ were already available.
\end{proof}

\section{Paths with free endpoints}
For $d\geq3$ set
\[
 Q(x)=\sum_{i=1}^d x_i^2+\sum_{i=1}^{d-1}\alpha_i x_ix_{i+1},
 \qquad \alpha_i\ne0.                                \tag{3}
\]
Every interval of the basis is orthogonal: its tridiagonal polar
matrix is nondegenerate in even dimension and has radical dimension
one in odd dimension. This follows successively from its first row;
equivalently the even determinant is the square of the product of
its alternate nonzero adjacent coefficients.

Choose $\lambda\ne0$ such that
\[
 H_\lambda=\begin{pmatrix}1&\lambda\\
                  \lambda&1+\lambda^2\end{pmatrix}
 \quad\hbox{has order }q+1.                          \tag{4}
\]
Such a choice always exists: take $\zeta$ of exact order $q+1$
in the norm-one torus of $\F_{q^2}^\times$ and set
$\lambda^2=\zeta+\zeta^{-1}\in F$.
The eigenvalues are $\zeta,\zeta^{-1}$, and are distinct.
The associated two-space is asingular, since a hyperbolic plane's
rotation group has order $q-1$.

\begin{theorem}\label{path}
Suppose (3) is nonsingular. If $d=3$, require at least one of its
two bonds to have the full order $q+1$ of (4).
If $d\geq4$, require every interior bond
$\alpha_2,\ldots,\alpha_{d-2}$ to have that order; the two
endpoint coefficients remain arbitrary nonzero elements.
Then its $d$ basis symmetries generate the full isometry group
and form a string C-group of rank $d$.
\end{theorem}
\begin{proof}
First prove generation of $\Sigma_U$ by the marked symmetries
of every interval $U$ of length at least three.
In length three use its full asingular two-space and
Lemma~\ref{three}; the third basis vector is nonradical because
its pairing with the middle vector is nonzero.

In length $b\geq4$, relabel the interval $v_1,\ldots,v_b$
and put $W=\gen{v_2,\ldots,v_b}$. By induction its marked
group is $\Sigma_W$. The full asingular line
$L_0=\gen{v_2,v_3}$ is available. Its conjugate
$L=L_0^{s_{v_1}}$ is asingular and not contained in $W$,
since $s_{v_1}(v_2)=v_2+\alpha_1v_1$.
The marked group contains $\Sigma_L$ by conjugation.
The argument of Lemma~\ref{saturation}, applied inside $U$ with
all symmetries still acting on the original ambient $V$, now gives
$\Sigma_U$. It uses only orthogonality of its three-spaces, so
also applies when the radical of this particular interval is singular.
With $U=V$, Lemma~\ref{support} gives the full isometry group.

For intersection, every two-generator submarking is a dihedral
C-group: its roots are independent, the symmetries are distinct,
and each has order two. In rank three, the intersection of the
end parabolics is supported on the middle one-space and hence is
exactly its symmetry subgroup by Lemma~\ref{support}.
For a longer interval, its two end parabolics are the full symmetry
groups of their spans. Their intersection is the symmetry group of
the middle span, again by Lemma~\ref{support}. In rank four this
middle span is the full-order interior two-space; in larger ranks
its generation was just proved. Apply the standard string C-group
end-parabolic criterion inductively. This criterion is the retained
McMullen--Schulte Proposition 2E16, also used in the accepted
L037 common contraction paper. Nonadjacent symmetries commute;
adjacent ones do not because their pairing is nonzero. All string
and full intersection requirements follow.
\end{proof}

\section{Both types and the odd-dimensional group}
\begin{theorem}\label{types}
For $q\geq4$, the groups $O^+(2m,q)$ and $O^-(2m,q)$ have
rank-$2m$ string C-markings by symmetries for every $m\geq1$.
The group $\Sp(2m,q)$ has one of rank $2m+1$ for every $m\geq1$.
Every adjacent product in these markings has odd order.
\end{theorem}
\begin{proof}
For dimension two the two orthogonal groups are the dihedral groups
of orders $2(q-1)$ and $2(q+1)$, generated by two symmetries whose
product generates the corresponding torus. Both torus orders are
odd and at least three.

For even $d=2m\geq4$, put every interior coefficient equal to
$\lambda$ of (4). Set the last coefficient to $\beta\ne0$,
and consider the odd suffix $W=\gen{v_2,\ldots,v_d}$.
Its radical generator has successive nonzero alternate coordinates
\[
                    z=v_2+v_4+\cdots+v_{d-2}
                              +(\lambda/\beta)v_d.
\]
Consequently $c=Q(z)=(m-1)+\lambda^2/\beta^2$.
If $m$ is even choose $\beta\ne\lambda$; if $m$ is odd any
nonzero $\beta$ works. Thus $c\ne0$.
Fix a nondegenerate complement $M$ to $\gen z$ in $W$.

Now let $\alpha=\alpha_1$ vary over $F^\times$.
The pairing $B(v_1,-)$ on $M$ is $\alpha$ times a fixed
functional, so there is a fixed $u\in M$ for which
$v'_1=v_1+\alpha u$ is perpendicular to $M$.
Write $Q(v'_1)=1+\alpha^2 h$, where $h$ is independent of
$\alpha$. Also $B(v'_1,z)=\alpha$.
The decomposition
$V=M\perp\gen{z,v'_1}$ therefore has Arf trace
\[
 \Tr(\operatorname{Arf} M+ch)+\Tr(c/\alpha^2).        \tag{5}
\]
Since $c\ne0$ and squaring is a bijection, the second term
takes both values zero and one as $\alpha$ ranges over
$F^\times$: a nonzero trace functional has $q/2$ points in
each fiber, and $q/2-1\geq1$ nonzero points in its zero fiber.
Both orthogonal types are thus attained. Theorem~\ref{path}
applies with these free endpoint choices.

For odd $d=2m+1$, take $\alpha_1=\cdots=\alpha_{d-2}=\lambda$
and $\alpha_{d-1}=\beta\ne0$. Its radical generator is
$v_1+v_3+\cdots+v_{d-2}+(\lambda/\beta)v_d$ and has value
$m+\lambda^2/\beta^2$. Choose $\beta\ne\lambda$ when $m$
is odd; otherwise any nonzero $\beta$ works.
The space is nonsingular with one-dimensional polar radical,
so Theorem~\ref{path} applies (including the rank-three case).
Its full isometry group is $\Sp(2m,q)$: the quotient map is
injective because the radical is nonsingular, and the lifted
transvections in Lemma~\ref{support} give surjectivity.

Finally an adjacent support is a nondegenerate two-space, and
the characteristic polynomial of its product is
$X^2+\alpha_i^2X+1$. It has distinct roots because
$\alpha_i\ne0$, lying in $F$ or its quadratic extension
as an inverse pair. Their orders divide $q-1$ or $q+1$,
hence are odd. On the perpendicular complement the product
is the identity, so its whole-space order is the same.
\end{proof}

\paragraph{Complete source interface.}
Theorem~\ref{types} gives the positive existence and exact group
claims of [BFL, Theorem 1.4 and Corollary 1.5], with a different,
fully specified type-selection proof. It does not assert that every
choice of source coefficients has nonsingular radical or generates
the entire isometry group when that radical is singular.
The three binary examples and higher-dimensional binary exclusion
are separate, fully acquired companion proofs with distinct reviews.
The source-interface and four-matrix companions, together with the
acquired small-dimension obstruction and centralizer papers, prove
the complete corrected low-dimensional spectrum.

\begin{thebibliography}{9}
\bibitem{BFL} P. A. Brooksbank, J. T. Ferrara and D. Leemans,
\emph{Orthogonal groups in characteristic 2 acting on polytopes
of high rank}, arXiv:1709.02219v2, 10 March 2018,
supplied twelve-page edition.
\bibitem{T} D. E. Taylor, \emph{The Geometry of the Classical Groups},
Heldermann, 1992, Theorems 11.39 and 11.43, printed pp.157 and 160.
Classical symmetry generation and quadratic forms remain retained inputs
with their full external proof costs.
\bibitem{LMO} D. Leemans, J. Moerenhout and E. O'Reilly-Regueiro,
\emph{Projective linear groups as automorphism groups of chiral
polytopes}, supplied twenty-five-page edition, Proposition 2.3, p.5;
the underlying Dickson classification remains an external input.
\bibitem{MS} P. McMullen and E. Schulte, \emph{Abstract Regular
Polytopes}, Cambridge, 2002, Proposition 2E16.
\end{thebibliography}
\end{document}
